\section{Model class}\label{sec:model}

\subsection{State}
A \tricl\ model has a finite set of \emph{entities}, each of one \emph{entity type}, and a finite set of \emph{relationship types}.
The state at time $t$ is a set $G(t)$ of \emph{links} $\ell = (e_1, r, e_3)$, each a directed edge of relationship type $r$ from entity $e_1$ to entity $e_3 \neq e_1$.
A relationship type may be declared \emph{symmetric} (e.g.\ \ttt{knows}) or have a named \emph{inverse} (\ttt{follows}/\ttt{is followed by}); then every link implies its mirrored link, and the two are established and terminated together.
A built-in identity relationship ``='' links every entity to itself and never changes; it lets rules refer to an existing link between the two entities, or to a link of one of them to a third entity of a given kind, with the same syntax as to two-paths.
Apart from the identity relationship there are no self-links, as in the original implementation: a self-link $(e_1, r, e_1)$ would turn every link $(e_1, r', e_3)$ into a two-path $e_1 \xrightarrow{r} e_1 \xrightarrow{r'} e_3$ of the same type as a two-path through a third entity, so that, e.g., an existing acquaintance would count as a common acquaintance; states of entities are encoded by links to auxiliary entities instead (Sec.~\ref{sec:expressiveness}).
The \emph{link type} of $\ell$ is the triple (type of $e_1$, $r$, type of $e_3$).

\subsection{Events and rates}
An \emph{event} establishes a link that does not exist or terminates one that exists.
For a candidate link $\ell = (e_1, r, e_3)$, an \emph{angle} is a two-path $e_1 \xrightarrow{r_{12}} e_2 \xrightarrow{r_{23}} e_3$ in $G(t)$ with $r_{12}, r_{23}$ including the identity relationship; its \emph{angle type} is $(r_{12}, \text{type of } e_2, r_{23})$.
For every event type $\tau$ (establish or terminate, link type), the configuration file specifies a basic attempt rate $a^0_\tau \ge 0$, basic success \emph{probability units} $b^0_\tau \in \R \cup \{\pm\infty\}$, and for any angle type $k$ an additional attempt rate $a_{\tau k}$ and additional probability units $b_{\tau k}$ contributed by each adjacent angle of that type.
With $n_k(\ell, t)$ the number of angles of type $k$ adjacent to $\ell$ at time $t$, the event happens at rate
\begin{equation}\label{eq:rate}
\begin{aligned}
  \lambda_\tau(\ell, t) &= A \cdot f(B), \\
  A &= a^0_\tau + \textstyle\sum_k a_{\tau k}\, n_k(\ell,t), \\
  B &= b^0_\tau + \textstyle\sum_k b_{\tau k}\, n_k(\ell,t),
\end{aligned}
\end{equation}
i.e.\ attempts occur at the rate $A$ and succeed with probability $f(B)$.
Coefficients may be expressions of named \emph{metaparameters}, which can be overridden on the command line and are the parameters estimated in Sec.~\ref{sec:inference}.
Infinite coefficients encode logic: a contribution $-\infty$ to $B$ makes the event impossible, and an infinite attempt rate with positive success probability makes it \emph{immediate}; all pending immediate events are performed, in a uniformly random order, before time advances.
Immediate events implement state machines: if an agent acquires a link \ttt{is infectious for} to the entity \ttt{covid-19}, the rule ``terminate \ttt{is susceptible to} immediately when the angle \ttt{is infectious for}, \ttt{covid-19}, = exists'' removes its previous state.

\subsection{The sigmoid}
The success probability is a sigmoid $f$ of the probability units with two \emph{tail indices} $\sigma_0, \sigma_1 \ge 0$ that set how fast it approaches $0$ and $1$:
\begin{equation}\label{eq:sigmoid}
\begin{aligned}
  f(B) &= \tfrac12 T_{\sigma_0}(-v) + \tfrac12 - \tfrac12 T_{\sigma_1}(v), \quad
  v = \frac{4B}{k(\sigma_0) + k(\sigma_1)}, \\
  T_\sigma(x) &= \bigl(1 + \sigma \ln(1 + e^{x})\bigr)^{-1/\sigma}, \quad T_0(x) = \frac{1}{1+e^{x}}, \\
  k(\sigma) &= (1 + \sigma\ln 2)^{-1-1/\sigma}, \quad k(0) = \tfrac12.
\end{aligned}
\end{equation}
$T_\sigma$ replaces the exponential tail of the logistic function by the $q$-exponential $(1 + \sigma y)^{-1/\sigma}$ of $y = \ln(1+e^{x})$ \citep{tsallis1988}.
For $\sigma > 0$ the tail decays like $|B|^{-1/\sigma}$, for $\sigma = 0$ exponentially.
Since $k(\sigma) = -2T'_\sigma(0)$, $f$ has slope $1$ at $B = 0$ for all tail indices, so that near zero, probability units are differences of probabilities. $f$ is strictly increasing, continuous in the tail indices, and equals $\expit(4B)$ for $\sigma_0 = \sigma_1 = 0$.
Heavy tails (the default is $\sigma_0 = \sigma_1 = 1$) keep rare events possible when many angles push the probability units far from zero.

\subsection{A model file}
The following excerpt is the dynamics of a model of triadic closure with triangle stabilization in three blocks of 100 entities (the \ttt{parameters\_3blocks} example):
\begin{lstlisting}[language=yaml]
relationship types:
  knows: symmetric
dynamics:
  [user, knows, user]:
    establish:
      attempt:
        basic: 0.001                     # spontaneous
        [~, knows, user, knows, ~]: 0.01 # per common
      success:                           # acquaintance
        basic: .inf                      # always succeed
    terminate:
      attempt:
        basic: 10.0                      # fast decay ...
      success:                           # ... hampered
        [~, knows, user, knows, ~]: -0.3 # by triangles
\end{lstlisting}
\noindent The angle type \ttt{[\~{}, knows, user, knows, \~{}]} counts common acquaintances. Further
sections of the file declare entity types and numbers, named entities, and initial links, which can be listed, read from csv files, or drawn from random graph models (Erd\H{o}s--R\'enyi, block models, random geometric graphs).

\subsection{Expressiveness}\label{sec:expressiveness}
Entity states are links to auxiliary ``hub'' entities: an agent is active if it has the link \ttt{is}~\ttt{active}.
A rule on the angle (agent, \ttt{knows}, agent, \ttt{is}, \ttt{active}) then depends on the number of active neighbors, which yields threshold models \citep{granovetter1978,wiedermann2020} (Fig.~\ref{fig:frames}) and contact processes; angles through ``='' make rates depend on an entity's own state, which yields compartmental epidemic models on adaptive networks \citep{gross2006}, including social distancing as termination of contacts with infectious agents.
Structural balance dynamics \citep{heider1946,antal2005} needs two relationship types \ttt{likes} and \ttt{dislikes} and four angle rules.
Well-mixed interactions are angles through a hub to which all entities are linked.
What is outside the class: rates that depend on structures beyond two-paths (e.g.\ on cycles of length four) and rates that depend on time or on the history of the process.
Rules on single links of one endpoint (``legs'', which would make rates depend on degrees) are part of the configuration syntax but rejected by the current version.
